\textbf{Čeření} odstraňuje drobné plynové bubliny ze skloviny. Přednáška rozlišuje:
\begin{itemize}
\item \textbf{Chemické čeření:} čeřiva, např. Na$_2$SO$_4$ či směs As$_2$O$_3$ a KNO$_3$, za vysoké teploty uvolňují plyny. Ty postupují k povrchu a strhávají s sebou bubliny z objemu skloviny.
\item \textbf{Fyzikální čeření:} únik bublin se usnadní snížením viskozity skloviny při vyšší teplotě; ve viskózní sklovině jinak bubliny stoupají pomalu.
\end{itemize}
Oba způsoby lze kombinovat.
